\documentclass{beamer}
\usepackage{amsfonts,amsmath}
\usepackage[portuguese]{babel}
\usepackage[utf8]{inputenc}
\usepackage[T1]{fontenc}
\usepackage{booktabs}
\usepackage{graphicx}
\usepackage{listings}
\usepackage{ragged2e}
\usepackage[normalem]{ulem}

\DeclareUnicodeCharacter{200B}{}
\usetheme{ufms}
\addtobeamertemplate{frame begin}{}{\justifying}

\usefonttheme[onlymath]{serif}

% Remover slides automaticos de sumario
\makeatletter
\let\beamer@writeslidentry@miniframeson=\beamer@writeslidentry
\def\beamer@writeslidentry@miniframesoff{\expandafter\beamer@ifempty\expandafter{\beamer@framestartpage}{}{\addtocontents{nav}{\protect\headcommand{\protect\beamer@partpages{\the\c@page}{\the\c@page}}}}}
\makeatother
\AtBeginSection[]{}

\titlebackground*{assets/background}

\newcommand{\hrefcol}[2]{\textcolor{cyan}{\href{#1}{#2}}}

\title{Graph-as-Code}
\subtitle{Origem, método e validações}
\course{PPGCC - FACOM/UFMS}
\author{\href{mailto:luiz.quirino@ufms.br}{Luiz F. P. \textbf{Quirino}}}
\IDnumber{}

\begin{document}
\maketitle

\begin{frame}{}
  \vspace{2em}
  \begin{center}
    \large
    Esta apresentação é baseada no artigo\\[0.5em]
    \textbf{``Actions Speak Louder than Prompts''}\\[0.3em]
    \normalsize
    Finkelshtein, Cucerzan, Jauhar \& White\\[0.3em]
    \emph{ICLR 2026}
  \end{center}

  \vspace{1.5em}
  \begin{center}
    \normalsize
    Os exemplos, dados e resultados apresentados aqui seguem\\
    o texto original. A narrativa foi reorganizada para fins didáticos.
  \end{center}
\end{frame}


% ============================================================
% 1. MOTIVACAO (2 frames)
% ============================================================
\section{Motivação}
\begin{frame}{O cenário}
  \textbf{LLMs se tornaram solucionadores de propósito geral.}\\
  {\small Com escala, modelos de linguagem adquirem capacidades que ninguém programou: raciocínio em cadeia, geração de código, uso de ferramentas. Funciona muito bem em texto puro.}

  \vspace{1em}
  \textbf{Grafos textuais são a fronteira natural.}\\
  {\small Redes sociais, citações, e-commerce --- entidades conectadas onde cada nó carrega texto. Se o LLM entende texto e o grafo codifica relações, combinar os dois promete classificação sem treino.}

  \vspace{1em}
  \textbf{O desafio: classificar com poucos rótulos.}\\
  {\small Rotular nós em grafos é caro. O cenário real é ter 5--10\% de labels e precisar classificar o restante. LLMs raciocinam zero-shot --- podem resolver isso?}
\end{frame}

\begin{frame}{O problema}

  \vspace{2em}
  \begin{block}{Síntese da pergunta central (Finkelshtein et al., ICLR 2026)}
    \emph{``Qual a melhor forma de um LLM interagir com dados de grafos?''}
  \end{block}

  \vspace{1.5em}
  \begin{center}
    \small
    O artigo testa 3 paradigmas de interação e apresenta que\\
    \textbf{deixar o LLM explorar o grafo gerando código} supera os demais.
  \end{center}
\end{frame}
% ============================================================
% 2. GENEALOGIA (6 frames)
% ============================================================
\section{Genealogia}
\begin{frame}{O ponto de partida}
  LLMs são bons em linguagem natural, mas ruins em:
  \begin{itemize}
    \item Cálculos (erram contas simples)
    \item Raciocínio em múltiplos passos (degradam com a profundidade)
    \item Dados estruturados (grafos, tabelas)
  \end{itemize}

 \vspace{2em}
  \begin{block}{Exemplo}
    \emph{``quantos vizinhos do nó 133 são classe 5?''}
  \end{block}

  \vspace{0.5em}
  Sabe-se que o LLM pode contar errado.
  A solução surgiu em etapas. Cada uma resolve a limitação da anterior.
\end{frame}

\begin{frame}[fragile]{2022: Chain-of-Thought -- CoT (Wei et al., NeurIPS)}
  \textbf{Insight:} se o LLM explicita o raciocínio passo a passo, erra menos.

  \vspace{0.5em}
  \textbf{Sem CoT:}
  \begin{lstlisting}[basicstyle=\ttfamily\scriptsize]
User:    Classe do no 133?
LLM:     Classe 3.                              (errado)
  \end{lstlisting}

  \vspace{0.3em}
  \textbf{Com CoT:}
  \begin{lstlisting}[basicstyle=\ttfamily\scriptsize]
User:    Classe do no 133?
LLM:     O titulo menciona Reinforcement Learning.
         4 dos 5 vizinhos pertencem a classe 5.
         Logo, classe 5.                        (correto)
  \end{lstlisting}

  \vspace{0.5em}
  \textbf{Limitação:} aritmética por inferência falha em problemas complexos.
\end{frame}

\begin{frame}[fragile]{2023: PAL -- Program-Aided Language Models (Gao et al., ICML)}
  \textbf{Insight:} em vez de raciocinar sobre números, gerar código que calcula.

  \vspace{0.3em}
  \begin{lstlisting}[language=Python, basicstyle=\ttfamily\scriptsize]
neighbors_labels = [5, 5, 5, 5, 1]
most_common = Counter(neighbors_labels).most_common(1)[0][0]
  \end{lstlisting}

  \vspace{0.3em}
  O LLM traduz o problema em programa. Python executa e retorna \texttt{5}.

  \vspace{0.5em}
  \textbf{Limitação:} o LLM precisa receber \emph{todos os dados} no prompt para gerar o código. Grafo grande estoura a janela de contexto.
\end{frame}

\begin{frame}[fragile]{2023: ReAct -- Reasoning + Acting (Yao et al., ICLR)}
  \textbf{Insight:} buscar dados sob demanda, agindo iterativamente.

  \vspace{0.3em}
  \begin{lstlisting}[basicstyle=\ttfamily\scriptsize]
Thought:  Preciso dos vizinhos do no 133.
Action:   get_neighbors(133)  ->  [134, 135, 707, 2048]

Thought:  Preciso do label de cada vizinho.
Action:   get_label(134)  ->  5
Action:   get_label(135)  ->  5
Action:   get_label(707)  ->  5
Action:   get_label(2048) ->  1

Thought:  Maioria classe 5.  Answer: 5
  \end{lstlisting}

  \vspace{0.3em}
  \textbf{Limitação:} 6 chamadas atômicas para uma contagem simples. Lento e \textit{relativamente} caro.
\end{frame}

\begin{frame}[fragile]{2024: CodeGraph (Cai et al., arXiv)}
  \textbf{Insight:} e se o LLM gerasse código para resolver problemas sobre grafos?

  \vspace{0.3em}
  \begin{lstlisting}[language=Python, basicstyle=\ttfamily\scriptsize]
def count_edges(adj_list):
    return sum(len(v) for v in adj_list.values()) // 2

count_edges(graph)
  \end{lstlisting}

  \vspace{0.3em}
  O LLM aprende com exemplos (few-shot), gera programas e executa via interpretador. Evita erros aritméticos e torna o raciocínio auditável.

  \vspace{0.5em}
  \textbf{Resultados:} melhoria de 1.3\% a 58.6\% em tarefas de raciocínio sobre grafos (GraphQA), dependendo da tarefa.

  \vspace{0.5em}
  \textbf{Limitação:} prova de conceito. Poucos datasets, sem comparação com outros paradigmas (Prompting, Tool-use), sem validação em escala.
\end{frame}

\begin{frame}{Genealogia \xout{completa} Lógica}
  \begin{center}
    \footnotesize
    \begin{tabular}{rcp{7.5cm}}
      \textbf{2022} & $\bullet$ & \textbf{Chain-of-Thought}: pensa passo a passo \\
      & $|$ & \emph{Limitação: cálculos de cabeça erram} \\[0.4em]
      \textbf{2023} & $\bullet$ & \textbf{PAL}: delega cálculos para código \\
      & $|$ & \emph{Limitação: precisa de todos os dados antecipadamente} \\[0.4em]
      \textbf{2023} & $\bullet$ & \textbf{ReAct}: busca dados iterativamente \\
      & $|$ & \emph{Limitação: ações atômicas, sem composição} \\[0.4em]
      \textbf{2024} & $\bullet$ & \textbf{CodeGraph}: código + grafos (prova de conceito) \\
      & $|$ & \emph{Limitação: poucos datasets, sem validação} \\[0.4em]
      \textbf{2026} & $\bullet$ & \textbf{Graph-as-Code}: síntese validada em escala \\
      & & 14 datasets $\cdot$ múltiplos modelos $\cdot$ node classification \\
    \end{tabular}
  \end{center}

  \vspace{0.5em}
  Cada passo resolve a limitação do anterior. Graph-as-Code (GaC) é a convergência natural.
\end{frame}

\begin{frame}[fragile]{2026: Graph-as-Code (Finkelshtein et al., ICLR)}
  \textbf{Síntese:} código (PAL) + iteração (ReAct) + grafos.

  \vspace{0.3em}
  \begin{lstlisting}[language=Python, basicstyle=\ttfamily\scriptsize]
df.loc[133, 'features']
df.loc[df.loc[133, 'neighbors'], 'label'].value_counts()
  \end{lstlisting}

  \vspace{0.3em}
  Passo 1: consulta o texto do nó alvo.\\
  Passo 2: conta os labels dos vizinhos numa expressão só.\\
  Passo 3: LLM combina os sinais e decide.

  \vspace{0.5em}
  \textbf{2 expressões} fazem o que ReAct levou 6 chamadas.\\
  A diferença com CodeGraph: GaC é \textbf{iterativo com feedback}. Cada expressão é informada pelo resultado da anterior. CodeGraph gera programa inteiro de uma vez sem ver resultados intermediários.\\
  Incorpora e estende a ideia do CodeGraph (Cai et al., 2024).
\end{frame}



% ============================================================
% 3. COMO FUNCIONA (2 frames)
% ============================================================
\section{Como funciona}
\begin{frame}{A ideia central}
  \textbf{O grafo inteiro vira uma tabela. (DataFrame)} Cada linha é um nó:

  \vspace{0.3em}
  \begin{table}
    \centering
    \footnotesize
    \begin{tabular}{clcl}
      \toprule
      \texttt{id} & \texttt{features} (o texto) & \texttt{label} (a classe) & \texttt{neighbors} (conexões) \\
      \midrule
      133 & ``A Reinforcement Learning...''  & \textbf{?} & [134, 135, 707, 2048] \\
      134 & ``Q-Learning for Robot...''      & 5   & [133, 135] \\
      135 & ``Policy Gradient Methods...''   & 5   & [133, 134, 707] \\
      707 & ``Temporal Difference...''       & 5   & [133, 135] \\
      2048 & ``Image Segmentation...''       & 1   & [133] \\
      \bottomrule
    \end{tabular}
  \end{table}

  \vspace{0.5em}
  {\small
  \textbf{Um nó:} um artigo, 3 atributos: features, neighbors, label.

  \vspace{0.2em}
  \textbf{Uma aresta:} citação entre artigos.

  \vspace{0.2em}
  \textbf{Labels visíveis:} poucos revelados como referência. O 133 é o alvo.
  }
\end{frame}

\begin{frame}{O mecanismo: perguntar, receber, decidir}

  O LLM funciona como um investigador com acesso a um terminal:

  \vspace{0.5em}
  \begin{center}
    \large
    \texttt{LLM pergunta $\rightarrow$ Python responde $\rightarrow$ LLM pergunta $\rightarrow$ ... $\rightarrow$ decide}
  \end{center}

  \vspace{0.5em}
  \begin{enumerate}
    \item LLM recebe: ``classifique o nó 133'' + estrutura da tabela
    \item Gera uma pergunta em forma de código (expressão pandas)
    \item Python executa e devolve a resposta como texto
    \item LLM lê a resposta, decide se precisa de mais informação
    \item Quando confiante: emite a classe final
  \end{enumerate}

  \vspace{0.5em}
  \textbf{Por que funciona:} o LLM escolhe \emph{o que} investigar (criatividade).\\
  Python garante que a resposta é exata (execução). Código gerado = explicação auditável.
\end{frame}


% ============================================================
% 4. EXEMPLO (2 frames)
% ============================================================

\begin{frame}[fragile]{Exemplo: classificando o nó 133, passo a passo}

  \textbf{Step 1 --- ``Sobre o que é este nó?''}
  \begin{lstlisting}[language=Python, basicstyle=\ttfamily\scriptsize]
df.loc[133, 'features']
  \end{lstlisting}
  \vspace{0.1em}
  $\rightarrow$ \texttt{"A Reinforcement Learning Approach to Job-Shop Scheduling"}

  \vspace{0.5em}
  \emph{O LLM agora sabe: o texto fala de Reinforcement Learning.}

  \vspace{0.8em}
  \textbf{Step 2 --- ``O que os vizinhos dizem?''}
  \begin{lstlisting}[language=Python, basicstyle=\ttfamily\scriptsize]
df.loc[df.loc[133, 'neighbors'], 'label'].value_counts()
  \end{lstlisting}
  \vspace{0.1em}
  $\rightarrow$ \texttt{\{5: 4, 1: 1\}} \quad (4 vizinhos são classe 5, 1 é classe 1)

  \vspace{0.5em}
  \emph{O LLM agora sabe: a vizinhança confirma --- maioria é classe 5.}
\end{frame}

\begin{frame}{Exemplo: classificando o nó 133, passo a passo}
  \framesubtitle{Step 3 --- decidir}

  O LLM combina os dois sinais:

  \vspace{0.5em}
  \begin{enumerate}
    \item \textbf{Texto} sugere Reinforcement Learning $\rightarrow$ hipótese: classe 5
    \item \textbf{Vizinhança} confirma: 4 de 5 vizinhos são classe 5
    \item Sinais convergem $\rightarrow$ confiança alta
  \end{enumerate}

  \vspace{0.3em}
  \begin{center}
    \large \texttt{Answer: 5} \quad (correto!)
  \end{center}

  \vspace{0.5em}
  \textbf{Por que isso é poderoso:}
  \begin{itemize}
    \item 2 perguntas bastaram, não precisou ver o grafo inteiro
    \item Texto e estrutura se complementam (um confirma o outro)
    \item Se o texto fosse ambíguo, a vizinhança desempataria
    \item Se vizinhos não tivessem label, o texto sustentaria sozinho
  \end{itemize}

  \vspace{0.3em}
  \emph{Adaptação flexível: usa o sinal mais informativo disponível.}
\end{frame}

% ============================================================
% 5. VALIDACAO (3 frames)
% ============================================================
\section*{Validação}
\begin{frame}{Validação: escopo do estudo original}
  Finkelshtein et al. testaram em:
  \begin{itemize}
    \item \textbf{14 datasets}: Cora, PubMed, ArXiv, Wiki-CS, Reddit, Cornell, Texas...
    \item \textbf{Múltiplos modelos}: o4-mini, GPT-5, DeepSeek R1, Phi-4, Qwen, Llama...
    \item \textbf{Foco}: node classification (+ shortest path como validação complementar)
  \end{itemize}

  \vspace{0.5em}
  \textbf{Baselines comparados:}
  \begin{table}
    \centering
    \footnotesize
    \begin{tabular}{ll}
      \toprule
      Categoria & Métodos \\
      \midrule
      GNNs supervisionadas & GCN, GAT, GraphSAGE (SOTA -- State of the Art) \\
      Semi-supervisionado & Label Propagation \\
      LLM zero-shot & Prompting, Tool-use, \textbf{Graph-as-Code} \\
      \bottomrule
    \end{tabular}
  \end{table}
\end{frame}

\begin{frame}{Resultados: eficiência de tokens}
  \textbf{Consumo de tokens por nó} (datasets de texto longo, Tabela 6 do paper):
  \begin{table}
    \centering
    \footnotesize
    \begin{tabular}{lrrr}
      \toprule
      Dataset & Prompting 2-hop & GraphTool & \textbf{GaC} \\
      \midrule
      Reddit & 42.100 & 16.100 & \textbf{14.900} \\
      Photo & 104.700 & 43.300 & \textbf{41.500} \\
      Wiki-CS & 1.157.000 (estoura) & 61.700 & \textbf{57.100} \\
      \bottomrule
    \end{tabular}
  \end{table}

  \vspace{0.5em}
  \textbf{Achado principal:} GaC usa 20$\times$ menos tokens que Prompting 2-hop em Wiki-CS. Em grafos densos, Prompting simplesmente nao cabe no contexto.

  \vspace{0.3em}
  {\footnotesize Nota: os resultados de accuracy variam por dataset e modelo. O paper reporta resultados especificos por combinacao, nao faixas gerais. Consultar Tabelas 2--5 do artigo para numeros exatos.}
\end{frame}

\begin{frame}{Ablação e achados principais}
  Finkelshtein remove componentes para medir a contribuição de cada sinal:

  \vspace{0.3em}
  \begin{itemize}
    \item Truncar texto: Prompting piora. \textbf{GaC compensa com estrutura.}
    \item Remover arestas: Prompting piora. \textbf{GaC compensa com texto.}
    \item Esconder labels: ambos pioram, mas GaC degrada menos.
  \end{itemize}

\end{frame}

\begin{frame}{Ablação e achados principais}

  \vspace{0.3em}
  {\small Tradução direta da introdução do artigo (Finkelshtein et al., ICLR 2026):}

  \vspace{0.5em}
  \begin{quote}
    \textbf{Achados.} Nossa avaliação produz percepções e diretrizes-chave para a aplicação de LLMs a dados de grafo:

    \vspace{0.3em}
    \begin{itemize}
      \item Graph-as-Code alcança o melhor desempenho geral, com ganhos especialmente grandes em \textbf{grafos de texto longo ou de alto grau}, onde o prompting rapidamente esgota o orçamento de tokens.

      \vspace{0.3em}
      \item Todos os modos de interação LLM--grafo são eficazes em \textbf{grafos heterofílicos}, desafiando a suposição comum de que eles colapsam sob baixa homofilia.

      \vspace{0.3em}
      \item Graph-as-Code é capaz de \textbf{deslocar flexivelmente} sua dependência entre estrutura, features ou rótulos, aproveitando o tipo de entrada mais informativo.
    \end{itemize}
  \end{quote}
\end{frame}

% ============================================================
% 6. POR QUE FUNCIONA + SHORTEST PATH (2 frames)
% ============================================================

\begin{frame}{Por que Graph-as-Code funciona}
  Três razões técnicas:

  \vspace{0.3em}
  \begin{enumerate}
    \item \textbf{Eficiência de tokens} --- consulta sob demanda, não serializa tudo\\
      {\footnotesize 57k tokens (GaC) vs. 1.157k (Prompting) no Wiki-CS}

    \vspace{0.3em}
    \item \textbf{Composicionalidade} --- operações complexas em 1 expressão\\
      {\footnotesize \texttt{.value\_counts()} = buscar + filtrar + contar}

    \vspace{0.3em}
    \item \textbf{Precisão algorítmica} --- código executa, não ``pensa''\\
      {\footnotesize MSE (Mean Squared Error) = 0.0 em shortest path (próximo slide)}
  \end{enumerate}

  \vspace{0.5em}
  \textbf{Adaptação flexível (Achado 9):} quando o texto é enganoso, GaC desloca para sinais estruturais sem reconfiguração.
\end{frame}

\begin{frame}[fragile]{Precisão algorítmica: shortest path}
  Tarefa: distância mínima entre dois nós.

  \vspace{0.3em}
  \textbf{Prompting/Tool-use:} LLM conta saltos de cabeça e erra em grafos grandes.

  \vspace{0.2em}
  \textbf{Graph-as-Code:} LLM gera código que resolve exatamente:
  \begin{lstlisting}[language=Python, basicstyle=\ttfamily\scriptsize]
G = nx.from_pandas_edgelist(
    df.explode('neighbors'), source='id', target='neighbors')
nx.shortest_path_length(G, source=42, target=133)
  \end{lstlisting}

  \vspace{0.3em}
  Resultado: MSE = 0.0 (erro zero) para GaC.

  \vspace{0.3em}
  \textbf{Princípio:} LLM decide \emph{o que} calcular (criatividade).\\
  Python garante que o cálculo é correto (execução).
\end{frame}

% ============================================================
% 7. COMPARACAO GNNs (1 frame)
% ============================================================
\section*{GaC vs GNNs}
\begin{frame}{GaC vs. GNNs (Graph Neural Networks) supervisionadas}
  \begin{table}
    \centering
    \footnotesize
    \begin{tabular}{p{2.8cm}p{4.5cm}p{4.5cm}}
      \toprule
      & \textbf{GNNs (GCN, GAT...)} & \textbf{Graph-as-Code} \\
      \midrule
      Treino & GPU + milhares de exemplos & Zero-shot (LLM congelado) \\
      Accuracy & Geralmente maior (SOTA) & Menor (trade-off por flexibilidade) \\
      Poucos rótulos & Sofre & Funciona (10\% basta) \\
      Explicabilidade & Caixa-preta & Código = explicação \\
      Novo domínio & Retreinar & Funciona direto \\
      Latência & ms & s--min por nó \\
      \bottomrule
    \end{tabular}
  \end{table}

  \vspace{0.3em}
  GaC não substitui GNNs. É alternativa para cenários com poucos rótulos, poucos recursos, ou necessidade de explicabilidade.
\end{frame}

% ============================================================
% 8. LIMITACOES (1 frame)
% ============================================================
\section*{Limitações}
\begin{frame}{Limitações e fronteiras}
  \textbf{GaC não tem vantagem quando:}
  \begin{itemize}
    \item Grafo esparso (vizinhança cabe no prompt)
    \item Nós isolados (sem vizinhança para explorar)
    \item Latência importa (40--160s por nó)
    \item Todos os labels disponíveis (GNNs dominam)
  \end{itemize}

\end{frame}

\begin{frame}{Limitações e fronteiras}

  \vspace{0.5em}
  \textbf{GaC é a melhor opção quando:}
  \begin{itemize}
    \item Grafo denso + poucos rótulos + texto longo
    \item Explicabilidade necessária (código auditável)
    \item Novo domínio/idioma sem dados de treino
    \item Custo importa
  \end{itemize}

  % \vspace{0.5em}
  % \textbf{Custo:} 50 nós $\approx$ \$0.30--\$2 (API) ou R\$0 (Ollama local).
\end{frame}

% ============================================================
% 9. RESUMO + REFERENCIAS (2 frames)
% ============================================================
\section*{Conclusão}
\begin{frame}{Resumo}
  \begin{center}
    \large
    \emph{Em vez de descrever o grafo inteiro para o LLM,\\
    deixamos ele explorar por conta própria, gerando código.}
  \end{center}

  \vspace{1em}
  \begin{enumerate}
    \item Grafo representado como DataFrame pandas
    \item LLM gera expressões para explorar vizinhança sob demanda
    \item Executor seguro avalia e devolve resultado
    \item LLM itera até decidir
  \end{enumerate}

  \vspace{0.5em}
  Funciona onde Prompting estoura (grafos densos) e Tool-use é lento (ações atômicas).
\end{frame}

\begin{frame}{Referências}
  \footnotesize
  \begin{itemize}
    \item Finkelshtein et al. (2026). ``Actions Speak Louder than Prompts''. ICLR.
    \item Gao et al. (2023). ``PAL: Program-Aided Language Models''. ICML.
    \item Yao et al. (2023). ``ReAct: Synergizing Reasoning and Acting in LLMs''. ICLR.
    \item Cai et al. (2024). ``CodeGraph: Enhancing Graph Reasoning of LLMs with Code''. arXiv:2408.13863.
    \item Wei et al. (2022). ``Chain-of-Thought Prompting Elicits Reasoning in LLMs''. NeurIPS.
    \item Wei et al. (2022). ``Emergent Abilities of Large Language Models''. TMLR.
  \end{itemize}
\end{frame}

\backmatter
\end{document}
